\begin{tikzpicture}[x=1cm,y=-.87cm,every node/.style={font=\fontsize{12.2}{15}\selectfont,text=Ink},arr/.style={-{Stealth[length=2mm]},line width=.8pt,draw=Classical}]\path[use as bounding box] (0,0) rectangle (7.45,6.2);
\draw[Rule,line width=.8pt] (0.8,1)--(0.8,4.6);
\draw[Rule,line width=.8pt] (2.0,1)--(2.0,4.6);
\draw[Rule,line width=.8pt] (3.2,1)--(3.2,4.6);
\draw[Rule,line width=.8pt] (4.3999999999999995,1)--(4.3999999999999995,4.6);
\draw[Rule,line width=.8pt] (5.6,1)--(5.6,4.6);
\draw[Rule,line width=.8pt] (0.8,1.0)--(5.6,1.0);
\draw[Rule,line width=.8pt] (0.8,2.2)--(5.6,2.2);
\draw[Rule,line width=.8pt] (0.8,3.4)--(5.6,3.4);
\draw[Rule,line width=.8pt] (0.8,4.6)--(5.6,4.6);
\fill[Quantum] (2.6,2.8) circle (3pt);\node[anchor=center,text=Ink,font=\fontsize{16}{15}\selectfont] at (6.35,2.8) {i};
\node[anchor=center,text=Ink,font=\fontsize{16}{15}\selectfont] at (2.6,5.2) {j};
\node[anchor=center,text=Ink,font=\fontsize{11}{15}\selectfont,text width=6.9cm,align=center] at (3.7,0.4) {La coordinata appartiene al centro della cella.};
\node[anchor=center,text=Ink,font=\fontsize{11}{15}\selectfont,text width=6.9cm,align=center] at (3.7,5.9) {Riga e colonna identificano una sola cella.};

\end{tikzpicture}